\documentclass[11pt,a4paper]{article}
\usepackage[margin=1in]{geometry}
\usepackage[T1]{fontenc}
\usepackage[utf8]{inputenc}
\usepackage{lmodern}
\usepackage{hyperref}
\usepackage{xcolor}
\setlength{\parskip}{0.7em}
\setlength{\parindent}{0pt}

\hypersetup{
  colorlinks=true,
  linkcolor=blue!50!black,
  urlcolor=blue!50!black
}

\title{UrbanFix / CivicSignal Platform Design Summary}
\author{Repository-derived design document}
\date{22 April 2026}

\begin{document}
\maketitle

\section{Purpose}
This document consolidates the current repository design into one LaTeX artifact. It is based on the top-level README, TODO list, the files in \texttt{design/}, and the implemented client and server code. It captures both the features already present in the codebase and the near-term product direction, including a new front-page AI chatbot.

\section{Repository Overview}
The repository is organized into three primary areas:
\begin{itemize}
  \item \texttt{client/}: Next.js 16 frontend using React 19 and Supabase client auth.
  \item \texttt{server/}: Express-based TypeScript API using Supabase as the data and auth backbone.
  \item \texttt{design/}: architecture, API, database, constraints, feature, and user-flow documents.
\end{itemize}

Important repository documents reviewed for this summary:
\begin{itemize}
  \item \texttt{README.md}
  \item \texttt{TODO.md}
  \item \texttt{client/README.md}
  \item \texttt{design/README.md}
  \item \texttt{design/features.md}
  \item \texttt{design/userflow.md}
  \item \texttt{design/constraints.md}
  \item \texttt{design/databaseFeatures.md}
  \item \texttt{design/dbDesign.md}
  \item \texttt{design/api-design.yml}
\end{itemize}

\section{Product Vision}
The platform is an anonymous civic issue reporting system for citizens, NGOs, government staff, and administrators. Citizens report local issues, attach images, discuss and validate posts, and browse a ranked public feed. Institution-side users receive richer triage detail, operational summaries, and moderation capabilities. AI is used as an assistive layer for spam rejection, enrichment, prioritization, and now conversational guidance through a planned front-page chatbot.

\section{Implemented and Planned Features}
\subsection{Citizen-Facing Features}
\begin{itemize}
  \item Supabase-backed authentication for registration, login, logout, and Google OAuth.
  \item Anonymous-by-default profile setup with public alias, language preference, and onboarding state.
  \item Public feed available on \texttt{/} and \texttt{/feed}.
  \item Search and filter controls for issue lookup by text and workflow status.
  \item Ranked issue cards with severity, status, raise count, comment count, follow count, and media preview.
  \item Issue detail pages with comment threads and engagement actions.
  \item Commenting, raising, following, and abuse reporting for authenticated users.
  \item Post creation flow with category selection, description entry, image upload, anonymity toggle, manual location, and auto-detected location.
  \item Metadata fallbacks for categories and areas when backend tables are unavailable.
  \item Guest browsing mode for read-only public access.
\end{itemize}

\subsection{Institution and Admin Features}
\begin{itemize}
  \item Role-aware institution dashboard for NGO staff, government staff, and admins.
  \item Summary metrics: total posts, unresolved posts, high-priority posts, and resolved posts.
  \item Severity breakdown and priority queue for operational triage.
  \item Role-scoped institution actions:
  \begin{itemize}
    \item NGO staff: update case status and internal notes within organization scope.
    \item Government staff: NGO capabilities plus public workflow state changes.
    \item Admins: global scope, report review, moderation queue visibility, and organization reassignment.
  \end{itemize}
  \item Institution-only post detail with exact location, moderation state, case tracking, and optional report history.
  \item Admin moderation queue for reported posts.
  \item Audit-safe public workflow history through \texttt{post\_status\_history}.
\end{itemize}

\subsection{Backend and Data Features}
\begin{itemize}
  \item Express route-object architecture with dynamic route loading.
  \item JWT validation against Supabase JWKS.
  \item Optional, required, and no-auth route protection modes.
  \item Dynamic rate-limiting by IP, authenticated user, or hybrid keying.
  \item CORS restrictions for local development and production.
  \item Supabase database access for posts, profiles, categories, areas, comments, raises, follows, reports, summaries, and institution case views.
  \item Supabase Storage support for user-scoped post image uploads.
  \item Fallback categories and fallback areas when lookup tables are empty or unavailable.
\end{itemize}

\subsection{AI and Automation Features}
\begin{itemize}
  \item Backend spam filtering using Groq with Llama 3.3 70B Versatile before a post is accepted.
  \item AI assessment model in the design for enrichment status, severity, complexity, translated text, hazard tags, and confidence.
  \item Priority-based ranking informed by AI severity, engagement, and freshness.
  \item Planned asynchronous enrichment pipeline for translation, severity analysis, and summary refreshes.
  \item Planned front-page AI chatbot for guided civic assistance.
\end{itemize}

\subsection{Front-Page AI Chatbot Addition}
The new chatbot should appear on the public home/feed page as a lightweight conversational assistant. Its role is not to replace formal reporting, moderation, or institution workflows. Instead, it should reduce friction for new users and improve navigation.

Recommended chatbot responsibilities:
\begin{itemize}
  \item Answer common user questions about how to report an issue.
  \item Explain categories, anonymity, raises, follows, comments, and institution access.
  \item Help users decide whether to create a new post, open an existing issue, or contact an institution workflow.
  \item Provide guided prompts such as ``What happened?'', ``Where is it located?'', and ``Does this affect safety, water, roads, or sanitation?''.
  \item Suggest the correct next page, such as \texttt{/register}, \texttt{/posts/new}, \texttt{/feed}, or \texttt{/institution}.
  \item Surface safety disclaimers: urgent emergencies should go to emergency channels, not only the app.
\end{itemize}

Recommended chatbot constraints:
\begin{itemize}
  \item It must not expose private moderation information, exact coordinates, or institution-only data.
  \item It should be advisory and navigational, not authoritative for legal or emergency decisions.
  \item It should not silently create posts; final post submission must stay explicit.
  \item It should reuse the same privacy model as the rest of the product.
  \item It should support multilingual prompting in the future, consistent with the translation helpers already planned.
\end{itemize}

\section{Architecture Summary}
\subsection{Frontend}
The client is a Next.js app using the App Router. It includes:
\begin{itemize}
  \item \texttt{AuthProvider} for session tracking and profile hydration.
  \item Feed, post creation, post detail, login, registration, logout, and institution pages.
  \item Shared UI components such as the app frame, auth screen, state block, and post card.
  \item Environment-based API access and Supabase browser client access.
\end{itemize}

\subsection{Backend}
The server is an Express app with:
\begin{itemize}
  \item route discovery via \texttt{routeHandler.ts}
  \item auth middleware based on route authorization level
  \item rate limiting per route policy
  \item centralized application error handling
  \item Supabase service-role access for data operations
  \item health endpoints including Groq reachability and configured-model checks
\end{itemize}

\subsection{Data Layer}
Supabase is used for:
\begin{itemize}
  \item authentication
  \item storage
  \item relational application tables
  \item future row-level security enforcement
\end{itemize}

The current design model includes profiles, organizations, areas, categories, posts, post media, AI assessments, comments, raises, follows, reports, status history, institution case views, and summary aggregates.

\section{Current API Surface}
The implemented route groups align well with the OpenAPI design:
\begin{itemize}
  \item \texttt{/health} and \texttt{/health/groq}
  \item \texttt{/me}
  \item \texttt{/profiles/register}
  \item \texttt{/feed}
  \item \texttt{/categories}
  \item \texttt{/areas}
  \item \texttt{/posts}
  \item \texttt{/posts/:postId}
  \item \texttt{/posts/:postId/comments}
  \item \texttt{/posts/:postId/raises}
  \item \texttt{/posts/:postId/follows}
  \item \texttt{/posts/:postId/reports}
  \item \texttt{/institution/posts/:postId}
  \item \texttt{/institution/summaries/overview}
  \item \texttt{/institution/summaries/areas/:areaId}
\end{itemize}

\section{Privacy, Safety, and Access Model}
\begin{itemize}
  \item Citizens are represented publicly through aliases rather than direct identity details.
  \item Exact coordinates are stored privately and intended only for institution-authorized views.
  \item Institution data access is role-based and organization-aware.
  \item Moderation state and abuse reports are hidden from normal public users.
  \item AI-derived severity affects ranking and triage, but moderation remains a separate concern.
  \item Spam detection rejects non-genuine posts before insertion.
\end{itemize}

\section{Database and Domain Model Summary}
The entity model described in the design folder includes:
\begin{itemize}
  \item user identity: auth users, profiles, organizations
  \item geography: areas
  \item civic taxonomy: categories
  \item core reporting: posts and media
  \item engagement: comments, raises, follows
  \item moderation: reports and moderation states
  \item triage and audit: status history and institution case views
  \item analytics: daily summaries and ranking projections
\end{itemize}

\section{Known Gaps and TODOs}
The repository TODO list currently identifies:
\begin{itemize}
  \item clean the README files and design directory
  \item update the README to reflect current changes
  \item add Docker support
  \item add integration tests
\end{itemize}

Additional implementation gaps observed from the repository:
\begin{itemize}
  \item the client README is still the default Next.js scaffold
  \item server tests are not implemented beyond a placeholder
  \item asynchronous enrichment workers are designed but not yet fully realized as a separate runtime
  \item the AI chatbot is planned but not yet implemented
\end{itemize}

\section{Recommended Next Steps}
\begin{enumerate}
  \item Implement the front-page AI chatbot as a guided assistant embedded in the feed hero or a persistent help drawer.
  \item Document the chatbot API contract, prompt boundaries, privacy rules, and fallback behavior.
  \item Add integration tests for auth, posting, engagement, institution workflows, and moderation.
  \item Replace remaining scaffold documentation with project-specific setup and architecture guidance.
  \item Introduce a worker abstraction for enrichment and summary refresh operations.
\end{enumerate}

\section{Conclusion}
This repository already contains a strong full-stack foundation for an anonymous civic reporting platform: authentication, posting, feed ranking, public discussion, institution triage, and moderation workflows are all represented either in code or in closely aligned design documentation. The new front-page AI chatbot fits naturally into this architecture as a guidance layer that improves discoverability, onboarding, and feature access while preserving the platform's privacy, workflow, and moderation boundaries.

\end{document}
